\documentclass[12pt,a4paper]{article}

% ---- Encoding & Font ----
\usepackage{iftex}
\ifPDFTeX
  \usepackage[utf8]{inputenc}
  \usepackage[T1]{fontenc}
\else
  \usepackage{fontspec}  % XeTeX/LuaTeX (tectonic): Unicode nativo (·, —, ¿, ª, §)
\fi
\usepackage[spanish]{babel}
\usepackage{csquotes}

% ---- Page layout ----
\usepackage[top=2.5cm, bottom=2.5cm, left=2.5cm, right=2.5cm]{geometry}
\usepackage{setspace}
\onehalfspacing

% ---- Math ----
\usepackage{amsmath, amssymb, amsthm}
\usepackage{mathtools}
\usepackage{physics}
\usepackage{esint}  % \oiint

% ---- Graphics ----
\usepackage{graphicx}
\usepackage{tikz}
\usepackage{float}
\usepackage{caption}

% ---- Tables ----
\usepackage{array}
\usepackage{booktabs}
\usepackage{tabularx}
\usepackage{colortbl}

% ---- Colours ----
\usepackage{xcolor}
\definecolor{notebg}{HTML}{FFF3CD}
\definecolor{noteborder}{HTML}{FFC107}
\definecolor{boxred}{HTML}{F8D7DA}
\definecolor{boxredborder}{HTML}{F5C6CB}

% ---- Boxes ----
\usepackage{mdframed}
\usepackage{tcolorbox}
\tcbuselibrary{most}

\newtcolorbox{keybox}[1][]{
  colback=blue!5,
  colframe=blue!70!black,
  arc=4pt,
  boxrule=1pt,
  fonttitle=\bfseries,
  title={#1}
}

\newtcolorbox{warnbox}[1][]{
  colback=boxred,
  colframe=boxredborder,
  arc=4pt,
  boxrule=1pt,
  fonttitle=\bfseries,
  title={#1}
}

\newtcolorbox{notebox}[1][]{
  colback=notebg,
  colframe=noteborder,
  arc=4pt,
  boxrule=1pt,
  fonttitle=\bfseries,
  title={#1}
}

% ---- Hyperlinks & TOC ----
\usepackage[hidelinks]{hyperref}
\usepackage{bookmark}

% ---- Enumerate ----
\usepackage{enumitem}

% ---- Miscellaneous ----
\usepackage{icomma}
\usepackage{siunitx}
\usepackage{textcomp}
\usepackage[bottom]{footmisc}

% ---- Title ----
\title{\textbf{Clase 12: Velocidad de Deriva, Ley de Ohm, \\
        Resistividad y Resistencia, y el Modelo de Drude}\\[0.3em]
        \large{Densidad de corriente y velocidad de deriva, conductividad,
              resistencia de conductores y conducción microscópica}}
\author{Transcripto y expandido --- Curso Física III (OpenFING)}
\date{}

\begin{document}

\maketitle
\thispagestyle{empty}
\newpage

\tableofcontents
\newpage

% ===================================================================
\section{Velocidad de desplazamiento (deriva)}
% ===================================================================

De la Clase 11: $I = \iint_S \mathbf{J}\cdot\hat{\mathbf{n}}\,dS$. Para una
superficie pequeña y plana de área $A$, despreciando la variación de
$\mathbf{J}$:
\[
I = J\,A\cos\theta,
\]
con $\theta$ el ángulo entre $\mathbf{J}$ y $\hat{\mathbf{n}}$. Si la
superficie es normal a $\mathbf{J}$: $I = J\,A$.

\subsection{Dirección y sentido de J}

La dirección de $\mathbf{J}$ es la del desplazamiento promedio de los
portadores (su velocidad media $\mathbf{v}_d$). El sentido:
\begin{itemize}[nosep]
  \item Portadores \textbf{positivos}: $\mathbf{J}$ en el sentido de
        $\mathbf{v}_d$.
  \item Portadores \textbf{negativos}: $\mathbf{J}$ opuesto a
        $\mathbf{v}_d$.
\end{itemize}
$\mathbf{v}_d$ es la \textbf{velocidad de desplazamiento} o de
\textbf{deriva}.

\subsection{Deriva vs.\ velocidad real; régimen estacionario}

La velocidad \textbf{real} de los electrones es muchísimo mayor que
$v_d$: se mueven rápido pero \textbf{en todas las direcciones}.
$\mathbf{v}_d$ es solo el pequeño arrastre \textbf{promedio}:
\begin{itemize}[nosep]
  \item Sin campo: movimiento aleatorio, promedio nulo.
  \item Con campo: se superpone un leve desplazamiento medio según el
        campo.
\end{itemize}

\begin{notebox}[Estacionario $\neq$ equilibrio; velocidad límite]
  Con corriente constante el sistema es \textbf{estacionario} pero no
  está en equilibrio (equilibrio sería $\mathbf{E}=0$ dentro). No hay
  aceleración constante sino \textbf{velocidad} constante: los choques
  con la red actúan como una viscosidad efectiva, y el sistema alcanza
  una velocidad límite (como un cuerpo en un medio viscoso).
\end{notebox}

% ===================================================================
\section{Relación entre J y la velocidad de deriva}
% ===================================================================

\subsection{Deducción con el cilindrito}

Cilindrito pequeño alineado con el campo, donde $\mathbf{E}$ y
$\mathbf{J}$ son uniformes. Sean $A$ el área de la base, $L = v_d\,\Delta
t$ el largo, y $n$ el número de portadores por unidad de volumen.

\begin{keybox}[Idea clave]
  En el intervalo $\Delta t$, los electrones que cruzan $A$ son
  \textbf{todos} los del cilindrito: avanzan a $v_d$, así que los del
  borde lejano recorren justo $L=v_d\,\Delta t$ y alcanzan a cruzar. (Es
  un enunciado en promedio, válido porque los portadores son
  muchísimos.)
\end{keybox}

Número de electrones $=n\,(A L)=n\,A\,v_d\,\Delta t$. Carga que cruza
(carga $-e$ por electrón):
\[
\Delta Q = (-e)\, n\, A\, v_d\, \Delta t
\;\Longrightarrow\;
I = \frac{\Delta Q}{\Delta t} = -e\, n\, A\, v_d.
\]

Dividiendo por $A$:

\begin{keybox}[Densidad de corriente y deriva]
  \[
  \boxed{J = -e\, n\, v_d}
  \]
\end{keybox}

\subsection{Forma vectorial general}

Con portadores de carga genérica $q$:
\[
\boxed{\mathbf{J} = q\, n\, \mathbf{v}_d}.
\]
Con varios tipos de portadores (p.\,ej.\ un electrolito):
\[
\mathbf{J} = \sum_i q_i\, n_i\, \mathbf{v}_{d,i}.
\]

% ===================================================================
\section{Estimación de la velocidad de deriva en el cobre}
% ===================================================================

\subsection{El cálculo}

Cable de cobre: $D = 1{,}8$~mm, $I = 1{,}3$~A. Densidad de corriente
(sección circular, $J$ uniforme):
\[
J = \frac{I}{\pi D^2/4} \approx 51\ \frac{\text{A}}{\text{cm}^2}.
\]

Densidad de portadores (un electrón por átomo de cobre):
\[
n = \frac{\rho_{\text{masa}}}{M_{\text{molar}}}\, N_A
  \approx \frac{8{,}96\times10^{3}}{63{,}5\times10^{-3}}\,(6{,}02\times10^{23})
  \approx 8{,}49\times10^{28}\ \frac{\text{e}^-}{\text{m}^3}.
\]

Velocidad de deriva (módulo):
\[
\boxed{v_d = \frac{J}{e\, n} \approx 14\ \frac{\text{cm}}{\text{hora}}}
\]
(otros metales, del mismo orden: algunos cm/hora).

\subsection{La analogía de la manguera y la señal eléctrica}

A $14$~cm/hora un electrón tardaría más de un día en llegar a la lámpara,
pero la luz se prende al instante. \textbf{Analogía de la manguera
llena:} al abrir la canilla, sale agua por el otro extremo de inmediato
aunque el agua avance lento — sale la que ya estaba del otro lado.

\begin{keybox}[A retener]
  La \textbf{señal eléctrica} viaja muchísimo más rápido que la velocidad
  de deriva.
\end{keybox}

\subsection{Corrección cuántica: el principio de Pauli}

\begin{warnbox}[La estimación está mal]
  $v_d$ real es unas \textbf{$\sim$100 veces mayor}. El error: no
  \textbf{todos} los electrones de conducción participan. Por el
  \textbf{principio de exclusión de Pauli}, casi todos los estados están
  ocupados y solo se mueven los electrones con estados libres al lado
  ($\sim$1 centésimo a temperatura ambiente). Entonces el $n$ efectivo es
  $\sim$100 veces menor y $v_d$, $\sim$100 veces mayor.
\end{warnbox}

\begin{notebox}[Dos ideas]
  (1) $v_d \ll$ velocidad real de los electrones (enorme pero errática,
  de promedio nulo sin campo). (2) La señal eléctrica viaja más rápido
  que ambas.
\end{notebox}

% ===================================================================
\section{Ley de Ohm (forma microscópica)}
% ===================================================================

\subsection{Conductividad y resistividad}

En un material \textbf{homogéneo e isótropo}, $\mathbf{J}$ debe ser
paralelo a $\mathbf{E}$ (no hay dirección preferencial). Se define la
\textbf{conductividad}:

\begin{keybox}[Conductividad y resistividad]
  \[
  \boxed{\mathbf{J} = \sigma\, \mathbf{E}}, \qquad
  \boxed{\rho = \frac{1}{\sigma}} \quad (\mathbf{E} = \rho\,\mathbf{J}).
  \]
\end{keybox}

$\sigma$ escalar: a igual $E$, mayor $\sigma$ da mayor $J$ (mide cuánto
conduce el material).

\begin{notebox}[Notación]
  $\sigma$ ya se usó para la densidad superficial de carga; el contexto
  distingue.
\end{notebox}

\subsection{Materiales óhmicos (lineales)}

A priori $\sigma$ podría depender de $E$ (problema de mecánica
estadística). Pero, como con los dieléctricos, muchos materiales son
\textbf{lineales}: $\sigma$ es, en muy buena aproximación,
\textbf{independiente de $E$} (constante del material). Entonces
$\mathbf{J}=\sigma\mathbf{E}$ es proporcionalidad: la \textbf{Ley de
Ohm}. Estos materiales se llaman \textbf{óhmicos} (metales, electrolitos,
buenos conductores).

\textbf{Unidades:} $[\sigma]=\text{Siemens/m}$ ($1\,\text{S}=1\,$A/V);
$[\rho]=\Omega\cdot\text{m}$ ($1\,\Omega=1\,$V/A). El Ohm se usa mucho; el
Siemens, casi nada.

\subsection{Tabla de resistividades}

\begin{center}
\begin{tabular}{lcc}
\toprule
\textbf{Material} & $\rho\ (\Omega\cdot\text{m})$ & \textbf{Tipo} \\
\midrule
Plata & $1{,}62\times10^{-8}$ & conductor \\
Cobre & $\approx 1{,}7\times10^{-8}$ & conductor \\
Aluminio & $2{,}75\times10^{-8}$ & conductor \\
Silicio puro & $2{,}5\times10^{3}$ & semiconductor \\
Silicio dopado (N) & $\ll$ puro & semiconductor \\
Agua destilada & $2{,}5\times10^{5}$ & aislante \\
Vidrio & $10^{10}$--$10^{14}$ & aislante \\
\bottomrule
\end{tabular}
\end{center}

\begin{itemize}[nosep]
  \item \textbf{Cobre vs.\ plata:} el cobre conduce un poco peor, pero se
        usa por costo-beneficio.
  \item \textbf{Semiconductores:} $\rho$ muy sensible a impurezas; al
        \textbf{dopar} (p.\,ej.\ fósforo) cae drásticamente (no ocurre en
        metales).
  \item \textbf{Conductor vs.\ aislante:} diferencia radical ($\sim
        10^{13}$).
  \item El vidrio no es realmente óhmico ($\rho$ depende del campo).
\end{itemize}

% ===================================================================
\section{Resistencia y ley de Ohm macroscópica}
% ===================================================================

$\mathbf{J}=\sigma\mathbf{E}$ es propiedad del \textbf{material};
$\Delta V=R\,I$, de un \textbf{cuerpo}.

\subsection{Conductor filiforme: $R=\rho L/A$}

Cable cilíndrico de área $A$, largo $L$, campo uniforme. Entonces
$|\Delta V| = E\,L$, y como es óhmico ($E=\rho J$) con $J=I/A$:
\[
|\Delta V| = \rho\,L\,\frac{I}{A} = \left(\frac{\rho\,L}{A}\right) I.
\]

\begin{keybox}[Resistencia de un filiforme]
  \[
  \boxed{R = \frac{\rho\, L}{A}}
  \]
\end{keybox}

\begin{itemize}[nosep]
  \item $\rho$ depende solo del material; $R$, del material \textbf{y} la
        forma.
  \item Doble largo $\Rightarrow$ doble $R$; más ancho $\Rightarrow$
        menos $R$ (analogía de la manguera).
  \item $R$ no depende de la intensidad.
\end{itemize}

\subsection{Caso general: la resistencia depende de los bornes}

La proporcionalidad $\Delta V = R\,I$ vale para \textbf{cualquier}
conductor óhmico; la fórmula $R=\rho L/A$, \textbf{solo} para
filiformes. En un cuerpo general, $R$ depende también de \textbf{dónde se
ponen los bornes} (bornes cercanos $\Rightarrow$ trayecto corto
$\Rightarrow$ $R$ pequeña).

\begin{keybox}[Resumen]
  \[
  \boxed{\Delta V = R\,I \ \text{(general, óhmicos)}}, \qquad
  \boxed{R = \frac{\rho L}{A} \ \text{(solo filiformes)}}
  \]
\end{keybox}

% ===================================================================
\section{Materiales óhmicos vs.\ no óhmicos: el diodo}
% ===================================================================

En la curva $I$ vs.\ $\Delta V$:
\begin{itemize}[nosep]
  \item \textbf{Óhmico:} recta por el origen (resistor/resistencia).
  \item \textbf{No óhmico:} cualquier curva que no sea recta por el
        origen.
\end{itemize}

Ejemplo no óhmico: el \textbf{diodo}. Curva asimétrica: para voltajes
positivos la corriente crece muy rápido; para negativos permanece casi
nula. Conduce fácil en un sentido y bloquea en el otro.

% ===================================================================
\section{Modelo microscópico de la conducción (Drude)}
% ===================================================================

Sección \textbf{6.4}. Modelo \textbf{clásico} de juguete (la descripción
rigurosa es cuántica) para entender por qué los metales son óhmicos.

\subsection{Velocidad real vs.\ velocidad de deriva}

Sin campo, el electrón se mueve caóticamente, chocando y arrancando en
direcciones aleatorias. Con campo, las trayectorias se corren
\textbf{levemente}. La velocidad real tiene módulo casi fijo y enorme: en
el cobre $v_e \approx 1{,}6\times10^{6}\ \text{m/s}$ (otros materiales,
$\sim 10^{6}$~m/s). La deriva es una fracción diminuta.

\subsection{Hipótesis del modelo}

\begin{enumerate}[nosep]
  \item \textbf{Mecánica clásica} (debería ser cuántica).
  \item \textbf{Movimiento libre entre choques:} solo acelerado por
        $\mathbf{E}$ (uniformemente acelerado).
  \item \textbf{Choques con pérdida total de memoria:} arranca en
        dirección aleatoria, misma rapidez, sin depender de la
        trayectoria previa.
  \item \textbf{Sin interacción electrón-electrón} (sorprendente, pero
        buena aproximación).
\end{enumerate}

Se define el \textbf{tiempo libre medio} $\tau$ (tiempo promedio entre
choques), que casi no depende de $E$ (el campo es diminuto frente al
movimiento térmico).

\begin{notebox}[Próxima clase]
  Con estas hipótesis se calcula la velocidad de deriva (movimiento
  uniformemente acelerado entre choques) y se deriva la conductividad.
\end{notebox}

\end{document}
